\section{Grothendieck's frameworks and the verified archival locators}\label{sec:archive}
\paragraph{G1: two letters, and two numbering conventions.}
The official Montpellier catalogue identifies cote~16 as a dossier concerning correspondences and Poincar\'e--Hodge algebras \citep{Groth16}. Its original PDF was retrieved from the catalogue's actual viewer link, and the relevant pages were rendered and visually inspected. The file has $56$ PDF pages, including the cover; on all inspected content pages the printed archive folio is one less than the one-based PDF page number. The source hash and retrieval details are recorded in the audit notes.

The sketch's attribution of PDF pages~44--50 to a single letter of 4 January 1967 is incorrect. Pages~44--47, printed folios~43--46, are that letter. Page~48, printed folio~47, starts a distinct Coates letter whose manuscript date reads \emph{6.1.1966}. Its semisimplicity discussion is on PDF page~49, printed folio~48, and its Weil-conjecture passage on PDF page~50, printed folio~49. The CSG draft transcriptions L10d and L11d agree with this division and with the second date as written. We do not emend that year on chronological grounds and do not claim a CSG transcription error. The correction is to the supplied sketch's attribution.

\begingroup\small
\begin{longtable}{@{}>{\raggedright\arraybackslash}p{25mm}>{\raggedright\arraybackslash}p{44mm}>{\raggedright\arraybackslash}p{76mm}@{}}
\toprule
Findings & Exact verified framework locator & Relationship, with its limit\\
\midrule
\endhead
R20, R24E--F, P0--P3, P2b, P5, SQ8--SQ9 & Cote~16, \S7, PDF pp.~14--17 / printed folios~13--16; \S8, PDF p.~18 / folio~17 & The integral exterior-algebra Lefschetz construction and its correspondence-ring language organize the projector and pairing. They do not assert bounded effective cycle representatives.\\
R21, R24A--B, M4 & FGA, expos\'e~221, Th\'eor\`emes~3.1--3.2, pp.~259--260; relative Chow refinement cited in R21 & Projective parameter spaces turn bounded cycle presentations into closed loci. A fixed-Hilbert-polynomial theorem alone must not be mistaken for a bounded-degree theorem.\\
R22, R23, P4 & SGA~1, Expos\'e XII, Th\'eor\`eme~5.1; rational weight-one Hodge/isogeny equivalence as cited in R22 & Finite monodromy covers are algebraic; Hecke correspondences act through the appropriate tensor stabilizer. The modern Shimura and Andr\'e--Oort inputs are identified separately.\\
M2--M10, MR1--MR3 & FGA, expos\'e~221, Th\'eor\`emes~3.1--3.2; SGA~6, X, Appendix~\S7, (7.13.10), pp.~578--579 & Coherent-sheaf and cycle-class framework. The actual analytic openness uses Grauert, \S7, Satz~3, pp.~62--63. The numerical regularity and effective Chern estimates are new derivations, not an archival corollary.\\
SQ1--SQ7 & Cote~16, \S8, PDF p.~18 / folio~17; SGA~1, XII, 5.1; FGA~221, 3.1--3.2 & Finite covers, algebraic correspondences and proper cycle loci organize the supplier transfer. Schoen and Milne provide the modern geometric and invariant-theoretic inputs.\\
P6 & SGA~4, XVI, Corollaire~2.2; SGA~6, X, Appendix~\S7, (7.13.10), pp.~578--579 & Smooth proper specialization and compatibility with rational cycle classes; the mixed-characteristic cycle parameter theorem is stated separately in P6.\\
P5, A3 & Cote~16, first Coates letter, PDF pp.~44--47 / folios~43--46; second letter, PDF pp.~48--50 / folios~47--49 & Algebraic Lefschetz inverses, index pairings and semisimplicity concern rational correspondences. Their numerical positivity does not control integral saturation or the size of effective positive and negative parts.\\
A1, A2 & The correspondence framework of cote~16; specific modern inputs Eskin--Oh, Theorem~1.2, and the conductor formula, Theorem~1.3 & There is no archival theorem of effective Hecke equidistribution or Colmez heights in these pages. Giving their modern exact locators is preferable to inventing an old one.\\
\bottomrule
\end{longtable}
\endgroup

On PDF page~14, printed folio~13, the manuscript explicitly begins with a free integral module of rank $2n$, an alternating form nondegenerate over $\Z$, its exterior algebra, wedge multiplication and contraction. It defines the integral ring generated by those two operators. The parenthetical integral nondegeneracy qualification is visible in the original. This is an integral linear-algebra construction. Interpreting it on the cohomology of a principally polarized abelian variety does not supply an effective cycle theorem.

On PDF page~45, printed folio~44, the first letter displays the critical-degree map $H^{2n-2}\to H^{2n+2}$ in its equivalence discussion. The second letter's semisimplicity passage uses the K\"unneth and index hypotheses. These structural questions are compatible with P5, but the inverse Gram matrix in P5 shows exactly where rational denominators remain even after the realization is a Tate sum.

No original cote~44 scan was used in this run. The sketch's references to that separate Jacobian programme are retained as source pointers only. No finding is attributed to a nonexistent checked proof in runs H7, H8, or H9: those run packages were not supplied. The mapping in \texttt{IMPLICATIONS.md} distinguishes the sketch's labels from independently inspected run statements.

\subsection*{Consequences for Hodge, Tate and the standard conjectures}
M1 gives the full Hodge conjecture on the fourfold family through the cited Markman corollary. P4 gives algebraicity of the Weil plane throughout the fixed split eightfold family. Together with the full-special-unitary Hodge-group hypothesis, the invariant-theory criterion in Severn's Theorem~11.2 yields the full Hodge conjecture for all powers of those members. P5 describes the resulting rational Tate summand in the category of homological motives. Its dual algebraic correspondences also pass to numerical motives; no claim about integral Chow projectors or rational equivalence is being substituted for the homological statement.

The supplied CM theorem is \citet[Theorem~1.1]{OpenAI26}, with the parallel human proof \citet[Theorem~I]{vanHoften26}. The supplied report \citet[Theorems~2.1, 3.1, and~3.5]{Milne26} records, respectively, the CM-Hodge input, the full Tate conjecture for abelian varieties over finite fields, and the Hodge standard conjecture for abelian varieties. These are cited background results under the prompt's convention. They are not newly derived from the eightfold calculation. Their positivity concerns algebraic classes modulo the relevant equivalence relation; A3 explains why it does not alone construct bounded effective representatives of an arbitrary rational Hodge tensor. The prime-specialization distinction is made in the preceding section.
